\section*{Capítulo \ref{ch:b2:blood-circulation} --- O sangue e o sistema circulatório}

\begin{solution}{exo:b2:blood-circulation:1}
\omterm{def:b2:blood-circulation:blood}{Plasma} \qty{55}{\%} (água, \qty{70}{g/L} de proteínas, sais,
nutrientes, gases, hormônios); células \qty{45}{\%}. \omterm{def:b2:blood-circulation:blood}{Hemácias} $5\times
10^{12}$/L, discos de \qty{7}{\micro m}, transporte de oxigênio;
leucócitos $7\times 10^{9}$/L, \qtyrange{10}{15}{\micro m}, imunidade;
\omterm{def:b2:blood-circulation:blood}{plaquetas} $3\times 10^{11}$/L, fragmentos de \qty{2}{\micro m},
hemostasia.
\end{solution}

\begin{solution}{exo:b2:blood-circulation:2}
Peixe: coração $\to$ brânquias $\to$ corpo $\to$ coração, um único
circuito, e o corpo recebe o \omterm{def:b2:blood-circulation:blood}{sangue} com a baixa pressão que resta após
as brânquias. Mamífero: coração direito $\to$ pulmões (\qty{25}{mmHg})
$\to$ coração esquerdo $\to$ corpo (\qty{120}{mmHg}) $\to$ coração
direito. O circuito duplo dá ao corpo uma alta pressão sem expor os
pulmões a ela, e permite que os dois circuitos sejam regulados
separadamente.
\end{solution}

\begin{solution}{exo:b2:blood-circulation:3}
\omterm{def:b2:blood-circulation:vessels}{Artéria}: parede espessa com camadas elásticas e músculo --- conduz o
\omterm{def:b2:blood-circulation:blood}{sangue} a alta pressão e suaviza o pulso. \omterm{def:b2:blood-circulation:vessels}{Arteríola}: parede formada
sobretudo por músculo liso --- fixa a resistência e distribui o fluxo.
\omterm{def:b2:blood-circulation:vessels}{Capilar}: uma célula endotelial de espessura --- trocas. \omterm{def:b2:blood-circulation:vessels}{Veia}: parede
fina e distensível, com válvulas --- devolve o \omterm{def:b2:blood-circulation:blood}{sangue} a baixa pressão e
armazena a maior parte dele.
\end{solution}

\begin{solution}{exo:b2:blood-circulation:4}
$70\times 72 = \qty{5}{L/min}$, \qty{7200}{L} por dia, contra
\qty{5}{L} de \omterm{def:b2:blood-circulation:blood}{sangue}: o volume passa pelo coração \num{1440} vezes por
dia. Nenhum órgão conseguiria produzir ou consumir sete toneladas de
\omterm{def:b2:blood-circulation:blood}{sangue} por dia; o mesmo \omterm{def:b2:blood-circulation:blood}{sangue} precisa voltar --- ele circula.
\end{solution}

\begin{solution}{exo:b2:blood-circulation:5}
$\Delta P = 8\eta LQ/\pi r^{4} = 8\times 3\times 10^{-3}\times
0.4\times 8.3\times 10^{-5}/(\pi\times 2.44\times 10^{-8}) =
\qty{10}{Pa}$, \qty{0.08}{mmHg}: a aorta não custa nada; a pressão é
gasta nas \omterm{def:b2:blood-circulation:vessels}{arteríolas}.
\end{solution}

\begin{solution}{exo:b2:blood-circulation:6}
$(15/12)^{4} = 2.4$: resistência 2.4 vezes maior, fluxo reduzido a
\qty{41}{\%}. $(15/20)^{4} = 0.32$: resistência reduzida a um terço,
fluxo 3.2 vezes maior.
\end{solution}

\begin{solution}{exo:b2:blood-circulation:7}
Aorta $83/3 = \qty{28}{cm/s}$; \omterm{def:b2:blood-circulation:vessels}{capilares} $83/3000 =
\qty{0.028}{cm/s}$; trânsito $0.1/0.028 = \qty{3.6}{s}$. No exercício,
o débito aumenta cinco vezes e a área \omterm{def:b2:blood-circulation:vessels}{capilar} aberta talvez três vezes,
de modo que o tempo de trânsito cai para cerca de \qty{2}{s} --- ainda
suficiente para que o oxigênio saia.
\end{solution}

\begin{solution}{exo:b2:blood-circulation:8}
Extremidade arterial $35 - 25 = \qty{+10}{mmHg}$ (filtração);
extremidade venosa $15
- 25 = \qty{-10}{mmHg}$ (reabsorção). Com $P_c = 28$ na extremidade
venosa, o saldo é $+3$: o líquido é filtrado ao longo de todo o
comprimento e nada é reabsorvido --- edema, os tornozelos inchados da
insuficiência cardíaca.
\end{solution}

\begin{solution}{exo:b2:blood-circulation:9}
$40/66000 = \qty{0.61}{mol/m^3}$; $\pi = 0.61\times 8.314\times 310 =
\qty{1570}{Pa} = \qty{11.8}{mmHg}$; com metade da albumina,
\qty{5.9}{mmHg}. O sódio atravessa livremente a parede \omterm{def:b2:blood-circulation:vessels}{capilar}, de modo
que sua concentração é a mesma dos dois lados e ele não exerce pressão
osmótica através dela.
\end{solution}

\begin{solution}{exo:b2:blood-circulation:10}
$RC = \qty{2}{s}$: $120e^{-0.4} = \qty{80}{mmHg}$. $C = 1$: $RC = 1$,
$120e^{-0.8} = \qty{54}{mmHg}$ --- e o mesmo volume sistólico eleva mais
a pressão sistólica numa aorta rígida. Uma \omterm{thm:b2:blood-circulation:windkessel}{pressão de pulso} larga
significa um pico mais alto contra o qual o coração precisa ejetar,
mais trabalho e mais oxigênio para o mesmo débito.
\end{solution}

\begin{solution}{exo:b2:blood-circulation:11}
Repouso: $90/83 = \qty{1.08}{mmHg.s/mL}$; exercício: $105/417 =
\qty{0.25}{mmHg.s/mL}$, uma queda de quatro vezes. Como $R \propto r^{-4}$,
$r$ aumentou por um fator $4^{1/4} = 1.41$: as \omterm{def:b2:blood-circulation:vessels}{arteríolas} se dilataram
em média \qty{41}{\%}.
\end{solution}

\begin{solution}{exo:b2:blood-circulation:12}
A continuidade fixa a velocidade em cada nível pela seção transversal
total, e a árvore é construída de modo que essa seção seja mil vezes a
da aorta nos \omterm{def:b2:blood-circulation:vessels}{capilares} e volte a ser pequena nas \omterm{def:b2:blood-circulation:vessels}{veias}; a \omterm{thm:b2:blood-circulation:poiseuille}{lei de
Poiseuille} coloca a resistência nas \omterm{def:b2:blood-circulation:vessels}{arteríolas}, onde o músculo pode
alterá-la, e deixa os vasos largos quase sem perdas. Uma \omterm{def:b2:blood-circulation:circuits}{circulação
aberta} não tem \omterm{def:b2:blood-circulation:vessels}{capilares} para desacelerar o \omterm{def:b2:blood-circulation:blood}{sangue} num lugar definido,
nem vasos nos quais fixar uma resistência, nem um meio de dirigir o
fluxo para um órgão --- ela só consegue agitar.
\end{solution}

\begin{solution}{pb:b2:blood-circulation:1}
\textbf{1.} $70\times 72 = \qty{5.0}{L/min}$; \qty{7300}{L} por dia.
\textbf{2.} $7300/5 \approx 1450$ vezes.
\textbf{3.} Um quarto de \qty{57}{g} a 72 batimentos por minuto:
\qty{14}{g} $\times 4320 = \qty{62}{kg}$ por hora --- mais ou menos o
peso de um homem.
\textbf{4.} Nenhum alimento conseguiria fornecer, e nenhum tecido
consumir, o peso de um homem em \omterm{def:b2:blood-circulation:blood}{sangue} a cada hora; a única saída é que
o mesmo \omterm{def:b2:blood-circulation:blood}{sangue} volte ao coração.
\textbf{5.} Um cordão amarrado em torno do braço faz inchar as \omterm{def:b2:blood-circulation:vessels}{veias}
abaixo dele, não as de cima, de modo que o \omterm{def:b2:blood-circulation:blood}{sangue} das \omterm{def:b2:blood-circulation:vessels}{veias} se move em
direção ao coração; empurrando o \omterm{def:b2:blood-circulation:blood}{sangue} de uma \omterm{def:b2:blood-circulation:vessels}{veia} em direção à mão,
ele para nas válvulas e não pode ser forçado para trás --- as válvulas
só permitem o fluxo em direção ao coração.
\textbf{6.} Os \omterm{def:b2:blood-circulation:vessels}{capilares} que unem as \omterm{def:b2:blood-circulation:vessels}{artérias} às \omterm{def:b2:blood-circulation:vessels}{veias}; Malpighi os viu
no pulmão de uma rã em 1661.
\textbf{7.} Área $\pi\times 1.25^{2} = \qty{4.9}{cm^2}$; $83/4.9 =
\qty{17}{cm/s}$.
\textbf{8.} $8\times 3\times 10^{-3}\times 0.4\times 8.3\times
10^{-5}/(\pi\times 2.44\times 10^{-8}) = \qty{10}{Pa}$, \qty{0.08}{mmHg}.
\textbf{9.} $R_{1} = 8\times 3\times 10^{-3}\times 10^{-3}/(\pi\times
5.06\times 10^{-20}) = \qty{1.5e14}{Pa.s/m^3}$; in parallel,
$1.5\times 10^{14}/3\times 10^{6} = \qty{5.0e7}{Pa.s/m^3}$.
\textbf{10.} $8.3\times 10^{-5}\times 5.0\times 10^{7} = \qty{4200}{Pa}
= \qty{31}{mmHg}$.
\textbf{11.} \omterm{def:b2:blood-circulation:vessels}{Capilares} abertos $10^{10}$, cada um com
$\pi(3\times 10^{-6})^{2}
= \qty{2.8e-11}{m^2}$: \qty{0.28}{m^2} $= \qty{2800}{cm^2}$; $v =
83/2800 = \qty{0.03}{cm/s}$; trânsito $0.1/0.03 = \qty{3.3}{s}$.
\textbf{12.} Resistência $\times(1/0.9)^{4} = 1.52$: \qty{47}{mmHg}; o
coração precisa elevar a pressão arterial em \qty{16}{mmHg}, ou o
débito cai um terço.
\textbf{13.} $+10$, $-10$ e $0$ mmHg.
\textbf{14.} Filtração $10\times 2 = \qty{20}{L}$ por dia; reabsorção
$8.5\times 2 = \qty{17}{L}$.
\textbf{15.} \qty{3}{L} de linfa por dia.
\textbf{16.} $\pi_c = \qty{12.5}{mmHg}$: saldo de $+22.5$ na
extremidade arterial e de $+2.5$ na extremidade venosa --- filtração em
toda parte, nenhuma reabsorção: edema.
\textbf{17.} Saldos de $+10$ e $+3$, média de $+6.5$: filtração ao longo
de todo o \omterm{def:b2:blood-circulation:vessels}{capilar}, uns \qty{26}{L} por dia; os linfáticos não conseguem
escoá-los e o líquido se acumula nos tecidos, primeiro nos pés.
\textbf{18.} $4\times 10^{10}\times 2\pi\times 3\times 10^{-6}\times
10^{-3} = \qty{750}{m^2}$; $t = x^{2}/2D = (2\times 10^{-5})^{2}/
(2\times 10^{-9}) = \qty{0.2}{s}$.
\textbf{19.} $(93 - 3)/83 = \qty{1.08}{mmHg.s/mL}$.
\textbf{20.} $RC = \qty{2.2}{s}$; $120e^{-0.8/2.2} = 120\times 0.69 =
\qty{83}{mmHg}$.
\textbf{21.} $RC = \qty{1.1}{s}$; $120e^{-0.73} = \qty{58}{mmHg}$: a
\omterm{thm:b2:blood-circulation:windkessel}{pressão de pulso} se alarga de 37 para \qty{62}{mmHg}.
\textbf{22.} $120e^{-0.3/2.2} = \qty{105}{mmHg}$: numa frequência
rápida, a pressão quase não cai entre os batimentos.
\textbf{23.} $C\Delta P = 2\times 20 = \qty{40}{mL}$ armazenados; os
outros \qty{30}{mL} escoam pelas \omterm{def:b2:blood-circulation:vessels}{arteríolas} durante a própria sístole.
\textbf{24.} O ventrículo que se contrai comprime seus próprios vasos
na sístole, de modo que o músculo cardíaco é irrigado na diástole, pela
pressão que a aorta, ao se retrair, mantém; uma aorta rígida que deixa
a pressão diastólica despencar priva o coração de \omterm{def:b2:blood-circulation:blood}{sangue} entre os
batimentos.
\textbf{25.} Débito de \qty{5}{L/min}; queda arteriolar de
\qty{31}{mmHg}; filtrado líquido de \qty{3}{L} por dia para a linfa;
diastólica de \qty{83}{mmHg} para $C = 2$ e de \qty{58}{mmHg} para
$C = 1$.
\end{solution}
